\documentclass[12pt,a4paper]{article}
\usepackage[utf8]{inputenc}
\usepackage[T1]{fontenc}
\usepackage[spanish]{babel}
\usepackage{geometry}
\usepackage{xcolor}
\usepackage{enumitem}

\geometry{a4paper, margin=1in}

\title{Reporte de Revisión de Pull Request \\ \large Módulo: Reconocimiento DNS (Grupo 1)}
\author{Prof. César Rodríguez \\ \small Curso: Seguridad Informática}
\date{\today}

\begin{document}

\maketitle

\section*{Estado del Pull Request: \textcolor{red}{Rechazado (Solicitud de Cambios)}}

Estimados estudiantes del Grupo 1: tras auditar su \textit{Pull Request} (PR), se ha determinado que actualmente \textbf{no puede ser integrado} a la rama principal de la Suite de Auditoría. Se han detectado violaciones críticas a las reglas de colaboración del repositorio y al contrato de datos.

\section*{Hallazgos y Correcciones Necesarias}

\subsection*{1. Violación Crítica: Modificación de Archivos Ajenos}
Han roto la \textbf{Regla de Oro} del trabajo colaborativo descrita en el archivo \texttt{README.md}:
\begin{itemize}[label=$\bullet$]
    \item \textbf{\texttt{discovery.py}:} Sobrescribieron el trabajo del Grupo 3, revirtiendo código que ya había sido aprobado.
    \item \textbf{\texttt{requirements.txt}:} Eliminaron dependencias estructurales críticas del proyecto (\texttt{jsonschema}, \texttt{dnspython}) y agregaron módulos nativos de Python (\texttt{datetime}, \texttt{ipaddress}, \texttt{subprocess}) lo cual causará errores en la instalación del entorno virtual.
\end{itemize}
\textbf{Acción requerida:} Deben descartar todos los cambios introducidos en estos dos archivos en su rama de trabajo.

\subsection*{2. Contrato de Datos (\texttt{schema\_resultados.json})}
El diccionario que retorna su función no cumple con el esquema oficial:
\begin{itemize}[label=$\bullet$]
    \item Cambiaron la clave obligatoria \texttt{'error\_message'} por \texttt{'error'}.
    \item Inventaron la clave \texttt{'querytype'}, la cual no existe en el esquema de validación.
    \item El campo \texttt{'status'} solo acepta \texttt{'success'} o \texttt{'error'}. Si un dominio no tiene registros, el estado debe ser \texttt{'success'} (la consulta fue exitosa) con las listas de datos vacías, nunca \texttt{'empty'}.
    \item En la clave \texttt{'estudiante'}, deben usar el identificador estándar \texttt{'E1'}, no su nombre completo.
\end{itemize}

\subsection*{3. Limpieza de Código (Código Muerto)}
Olvidaron eliminar la instrucción \texttt{raise NotImplementedError(...)} al inicio de la función. Al estar posicionada antes del bloque \texttt{try-except}, el programa lanzará la excepción de inmediato y el resto de la lógica que programaron jamás se ejecutará.

\subsection*{4. Resiliencia (Puntos Positivos)}
\textcolor{green!70!black}{\textbf{Aprobado parcialmente:}} El manejo de excepciones específicas de la librería \texttt{dns.resolver} (tales como \texttt{NXDOMAIN}, \texttt{NoNameservers}, etc.) está excelentemente estructurado. Solo deben asegurarse de ajustar las variables dentro de los bloques \texttt{except} para que respeten la clave \texttt{'error\_message'}.

\section*{Instrucciones para la Corrección}
No es necesario que abran un nuevo Pull Request. Pueden arreglar su rama actual ejecutando los siguientes comandos en su terminal local. 

Primero, restauren los archivos que no les corresponden para que vuelvan a su estado original (el de la rama principal del curso):
\begin{verbatim}
git checkout origin/main -- modulos/discovery.py requirements.txt
\end{verbatim}

Luego, apliquen las correcciones a su código en \texttt{modulos/dns\_recon.py}. Una vez que prueben que todo funciona, suban los cambios a su rama:
\begin{verbatim}
git add modulos/dns_recon.py modulos/discovery.py requirements.txt
git commit -m 'Correcciones de PR: contrato de datos y limpieza de archivos ajenos'
git push origin <nombre-de-su-rama>
\end{verbatim}

\vspace{1cm}
\noindent Quedo a la espera de las correcciones para proceder con la aprobación y el pase a producción.

\end{document}